\section{Introduction}

The minimum-variance portfolio is the solution to a quadratic programme whose
inputs are a covariance estimate and a set of weight constraints. In the large-$n$
regime of institutional equity portfolios ($n/T \approx 0.4$ for a five-year daily
window), the sample covariance $\hat\Sigma = X^\top X/T$ is nearly singular,
with condition number $\kappa(\hat\Sigma) \approx 106{,}000$ on S\&P~500 data.
This ill-conditioning makes the portfolio optimisation problem numerically fragile.

Random matrix theory (RMT) attributes most of this ill-conditioning to noise:
the Marchenko--Pastur (MP) law identifies eigenvalues below the bulk edge
$\hat\lambda_+$ as uninformative, consistent with noise, while eigenvalues above
$\hat\lambda_+$ carry genuine factor structure. The ``constant residual eigenvalue''
(CRE) cleaning method~\cite{laloux1999,bun2017,bouchaud2003} replaces the noise
eigenvalues with their average $\bar\lambda$ -- preserving the trace of
$\hat\Sigma$ -- and produces the estimator
\[
  T_0^{\mathrm{RMT}} = \bar\lambda I + U_k(\Lambda_k - \bar\lambda I)U_k^\top,
\]
where $(U_k, \Lambda_k)$ are the $k$ eigenpairs above $\hat\lambda_+$.
On S\&P~500 data, $k = 23$ and $\kappa(T_0^{\mathrm{RMT}}) \approx 343$ --
roughly $300$ times better conditioned than $\hat\Sigma$.

Implementations of CRE cleaning in the public domain typically form
$T_0^{\mathrm{RMT}}$ as a dense $n \times n$ matrix and pass it to a
general-purpose solver (CVXPY, interior-point, etc.); see e.g.\ the reference
implementations accompanying~\cite{lopezdeprado2018}. This paper exploits the
low-rank-plus-identity structure
$T_0^{\mathrm{RMT}} = \bar\lambda I + U_k \Delta_k U_k^\top$, which renders
the matrix inverse available in closed form via the Woodbury identity, without
assembling either the $n \times n$ covariance or the $n_a \times n_a$ active-set
restriction. The result is a direct active-set portfolio solver that is
substantially faster than both general-purpose interior-point solvers and the
KKT-Cholesky approach of assembling the restricted system and factoring.
Exploiting factor structure in portfolio quadratic programmes via the
Sherman--Morrison--Woodbury identity goes back at least to
Perold~\cite{perold1984}; the contribution here is the combination with
RMT-determined structure, randomised preprocessing, and an end-to-end
complexity and accuracy account.

\paragraph{Contributions}
This paper makes four contributions, each self-contained relative to the
companion paper~\cite{fast_minimum_variance}.
The companion paper uses a different estimator structure and a different inner
solver; the Woodbury treatment, the crossover analysis, the rSVD preprocessing
treatment, and the total pipeline complexity accounting developed here are absent
from it.

\emph{First}, we give a complete treatment of the randomised SVD preprocessing
step (Section~\ref{sec:preprocessing}). The companion paper uses a dense
eigendecomposition of $\hat\Sigma$ for simplicity; at large $n$ this
costs $O(n^2 T + n^3)$ and re-introduces the $n^2$ storage the matrix-free
approach aims to avoid. A randomised truncated SVD applied directly to $X$
computes the top-$k$ eigenpairs of $\hat\Sigma$ in $O(nTk)$ time and $O(nk)$
storage; the measured portfolio impact of the approximation is below three
basis points in any weight. For daily rolling windows we additionally
implement an incremental rank-two update of the tracked eigenpairs that
avoids re-sketching altogether.

\emph{Second}, we derive the exact crossover condition for Woodbury vs.\
KKT-Cholesky (Proposition~\ref{prop:crossover}): the Woodbury inner solve
requires fewer flops per outer step precisely when $k < n_a$, a condition
satisfied throughout the long-only equity regime
($k/n_a$ between roughly $0.2$ and $0.7$ in our experiments). We also account
for the total pipeline cost over a sweep of $S$ related solves
(Proposition~\ref{prop:pipeline}), under which the $O(nTk)$ preprocessing
amortises to $O(nTk/S)$ per solve.

\emph{Third}, we provide a direct numerical comparison of Woodbury against
KKT-Cholesky, a tuned interior-point baseline (Clarabel called natively on
pre-assembled data), and CVXPY on the same RMT problem
(Section~\ref{sec:results}), together with optimality cross-checks of the
returned weights, an exact critical-line sweep built on the same Woodbury
kernel that traces every frontier breakpoint (Section~\ref{ssec:cla}), a
$k$-sensitivity audit showing that the MP threshold choice
is not a sensitive hyperparameter (Section~\ref{sec:ksensitivity}), and an
out-of-sample backtest on two universes -- S\&P~500 and FTSE~100 -- showing
that the cleaned estimator carries no measurable statistical penalty relative
to standard shrinkage (Section~\ref{ssec:oos}).

\emph{Fourth}, we delineate the scope of the structural trick
(Remarks~\ref{rem:rie} and~\ref{rem:ridge}): ridge terms, diagonal loadings,
externally specified factor models, and linear equality constraints preserve
the Woodbury structure, whereas statistically richer cleanings -- nonlinear
shrinkage and general rotationally invariant estimators -- do not.

\paragraph{Relation to existing work}
Laloux et al.~\cite{laloux1999} and Plerou et al.~\cite{plerou1999} introduced
the Marchenko--Pastur threshold to financial covariance estimation; Bouchaud and
Potters~\cite{bouchaud2003} develop the theory in detail.
Bun, Bouchaud and Potters~\cite{bun2017} give a comprehensive treatment of CRE
cleaning and of the statistically superior rotationally invariant estimators
built on the eigenvector overlaps of Ledoit and
P\'ech\'e~\cite{ledoitpeche2011}; nonlinear shrinkage is developed by Ledoit
and Wolf~\cite{ledoit2012,ledoit2020}. Those estimators modify the full
spectrum and are not Woodbury-compatible; this paper works with CRE precisely
because clipping is the member of the family with exploitable algebraic
structure (Remark~\ref{rem:rie}).
L\'opez de Prado~\cite{lopezdeprado2018} provides accessible CRE
implementations for practitioners, built on dense covariance assembly and
general-purpose solvers.

On the optimisation side, Perold~\cite{perold1984} exploited
diagonal-plus-low-rank factor structure inside a large-scale portfolio QP
solver four decades ago, and the critical line algorithm of
Markowitz~\cite{markowitz1956} (see~\cite{niedermayer2010,cvxcla} for modern
treatments) traces the exact long-only frontier by active-set pivoting.
Relative to this literature our contribution is the pairing of the
\emph{RMT-determined} low-rank structure with the active-set solve, the
randomised-SVD preprocessing with an explicit total-pipeline complexity
accounting, and the measured end-to-end comparison
(Remark~\ref{rem:cla} discusses the relation to CLA).
The randomised SVD itself is due to Halko, Martinsson and
Tropp~\cite{halko2011} and has been applied widely to PCA and large-scale
covariance estimation.

This paper is self-contained: Algorithm~\ref{alg:woodbury} gives the complete
primal-dual active-set solver including the Woodbury inner solve, and
Section~\ref{sec:preprocessing} develops the rSVD preprocessing independently.
A companion paper~\cite{fast_minimum_variance} addresses the same portfolio
problem with an iterative inner solver and a different covariance structure.

\paragraph{Outline}
Section~\ref{sec:problem} states the problem. Section~\ref{sec:rmt} constructs
the cleaned estimator. Section~\ref{sec:algorithm} develops the Woodbury portfolio
algorithm and the crossover analysis. Section~\ref{sec:preprocessing} covers
randomised SVD preprocessing. Section~\ref{sec:ksensitivity} analyses sensitivity
to the threshold $k$. Section~\ref{sec:results} reports timing, scaling, and
out-of-sample results. Section~\ref{sec:conclusions} concludes.
